\documentclass[11pt,a4paper]{article}

\usepackage[a4paper,margin=1in]{geometry}
\usepackage{fontspec}
\IfFontExistsTF{Times New Roman}
  {\setmainfont{Times New Roman}}
  {\IfFontExistsTF{TeX Gyre Termes}{\setmainfont{TeX Gyre Termes}}{\setmainfont{Times}}}
\usepackage{newtxmath}
\usepackage{amsmath}
\usepackage{booktabs}
\usepackage{longtable}
\usepackage{array}
\usepackage{enumitem}
\usepackage{xcolor}
\usepackage{listings}
\usepackage{graphicx}
\usepackage{float}
\usepackage[hidelinks]{hyperref}
\usepackage{xurl}
\usepackage{fancyhdr}
\usepackage{microtype}

\setlength{\parindent}{0pt}
\setlength{\parskip}{0.65em}
\setlength{\emergencystretch}{3em}
\setlist[itemize]{leftmargin=1.4em,itemsep=0.2em,topsep=0.2em}
\setlist[enumerate]{leftmargin=1.6em,itemsep=0.2em,topsep=0.2em}

\pagestyle{fancy}
\fancyhf{}
\lhead{\small Ticket 11 Study Report}
\rhead{\small Grid and Domain Sensitivity}
\cfoot{\thepage}
\renewcommand{\headrulewidth}{0.3pt}

\lstset{
  basicstyle=\ttfamily\small,
  breaklines=true,
  frame=single,
  columns=fullflexible,
  keepspaces=true,
  rulecolor=\color{black!30},
  xleftmargin=0.5em,
  xrightmargin=0.5em
}

\newcommand{\code}[1]{\texttt{#1}}
\newcommand{\file}[1]{\path{#1}}
\newcommand{\func}[1]{\path{#1}}
\newcolumntype{L}[1]{>{\raggedright\arraybackslash}p{#1}}
\newcolumntype{P}[1]{>{\raggedright\arraybackslash\ttfamily\footnotesize}p{#1}}

\title{\textbf{Ticket 11 Study Report: Grid and Domain Sensitivity}\\
\large SURF2026 American Option Risk Surfaces}
\author{Codex-assisted implementation report}
\date{June 16, 2026}

\begin{document}
\maketitle

\begin{abstract}
Ticket 11 studies whether representative American CN/PSOR outputs are stable when the grid
resolution and finite spot-domain cutoff change. It uses the baseline Ticket 04 solver, not the
Ticket 10A Rannacher variant. The study compares selected prices, continuation-premium boundary
curves, LCP diagnostics, PSOR iterations, and finite-difference Gamma diagnostics for an American
put and a dividend-paying American call. The evidence is encouraging for this preliminary ladder:
obstacle violations remain zero, complementarity diagnostics remain near solver tolerance, and
the proportional domain extension from \(S_{\max}=4\) to \(6\) does not change selected target
region values. The report is intentionally cautious because three grid levels are not a formal
convergence proof.
\end{abstract}

\tableofcontents
\newpage

\section{Role of Ticket 11 in the Solver Project}

Earlier tickets built the validation ladder: European closed-form utilities, grids and operators,
European CN validation, American CN/PSOR, American put and call validation, LCP diagnostics,
boundary extraction, Delta/Gamma diagnostics, and a Rannacher smoothing comparison. Ticket 10A
concluded \path{KEEP_BASELINE_FOR_NOW}, so Ticket 11 uses the baseline CN/PSOR solver.

Ticket 11 answers a different question: do representative prices, boundaries, and diagnostics
remain stable when we change \(M\), \(N\), and \(S_{\max}\)? This is required before stress maps,
dataset construction, or neural surrogate modelling.

\section{Theory and Methodology Connection}

The documented source basis is internal: the planning report, literature map, formulation note,
solver validation plan, and prior Ticket 01--10A reports. The solver validation plan explicitly
places grid refinement and domain sensitivity before stress maps and neural work. No BibTeX file
is created, and no external bibliographic details are invented.

\section{Grid Sensitivity in Plain Language}

A finite-difference solver replaces a continuous PDE problem with values on a grid. A coarse grid
can miss important behavior near the strike or exercise boundary. Grid sensitivity checks whether
the result changes less as the grid is refined.

In this ticket:
\[
M=N\in\{80,120,180\},\qquad S_{\max}=4.
\]
The finest grid, \(M=N=180\), is used as the numerical reference for selected comparisons.

\section{Domain Cutoff Sensitivity in Plain Language}

The Black--Scholes spot domain is theoretically unbounded, but the code solves on
\([0,S_{\max}]\). If \(S_{\max}\) is too small, the artificial upper boundary can distort prices
or boundaries inside the region we care about.

Ticket 11 tests:
\[
S_{\max}\in\{4,5,6\}.
\]
The domain cases use proportional \(M\) values:
\[
(S_{\max},M,N)\in\{(4,120,120),(5,150,120),(6,180,120)\}.
\]
This keeps \(\Delta S\approx 1/30\) comparable while changing the domain cutoff.

\section{Why Ticket 11 Follows Ticket 10A}

Rannacher smoothing is useful only if it is clearly reported and improves diagnostics without
distorting prices or boundaries. Ticket 10A did not find a clear Gamma benefit, so Ticket 11 does
not silently switch methods. This keeps the sensitivity study tied to the already validated
baseline solver.

\section{Sensitivity Design and Chosen Configurations}

\begin{longtable}{P{0.24\textwidth}P{0.18\textwidth}L{0.40\textwidth}}
\toprule
Study & Cases & Purpose \\
\midrule
\endhead
Grid sensitivity & \(M=N=80,120,180\), \(S_{\max}=4\) & Check refinement behavior while holding the domain fixed. \\
Domain sensitivity & \(S_{\max}=4,5,6\), \(M=120,150,180\), \(N=120\) & Check upper-boundary influence while keeping \(\Delta S\) comparable. \\
Option families & American put and dividend-paying American call & Reuse the representative cases from Tickets 05 and 07. \\
\bottomrule
\end{longtable}

\section{Reference Configuration Choice}

The reference is not an analytical truth. It is simply the finest or largest-domain configuration
available inside Ticket 11:
\begin{itemize}
  \item grid reference: \(M=N=180\), \(S_{\max}=4\);
  \item domain reference: \(S_{\max}=6\), \(M=180\), \(N=120\).
\end{itemize}

Every selected spot row reports actual moneyness. This matters because nearest-node comparisons
can compare \(S/K=0.5\) on one grid with \(S/K=0.4888889\) on another grid.

\section{Price Stability Results}

\begin{longtable}{P{0.27\textwidth}P{0.15\textwidth}P{0.14\textwidth}P{0.16\textwidth}P{0.16\textwidth}}
\toprule
Case & \(M,N\) or \(S_{\max}\) & Reference & Max selected diff & RMSE selected diff \\
\midrule
\endhead
\path{american_put_grid_80} & 80 & grid 180 & 0.0111111 & 0.00454879 \\
\path{american_put_grid_120} & 120 & grid 180 & 0.0111111 & 0.00453795 \\
\path{american_put_grid_180} & 180 & grid 180 & 0 & 0 \\
\path{dividend_call_grid_80} & 80 & grid 180 & 0.0111111 & 0.00454255 \\
\path{dividend_call_grid_120} & 120 & grid 180 & 0.0111111 & 0.00453671 \\
\path{dividend_call_grid_180} & 180 & grid 180 & 0 & 0 \\
\path{american_put_domain_smax4} & 4 & domain 6 & 0 & 0 \\
\path{dividend_call_domain_smax4} & 4 & domain 6 & 0 & 0 \\
\bottomrule
\end{longtable}

The largest grid selected-spot difference is driven by nearest-node matching at \(S/K=0.5\). For
example, the \(M=180\) grid's nearest spot to \(0.5K\) is \(0.4888889K\), while the coarser grids
hit \(0.5K\) exactly. This is a useful warning for beginner researchers: a selected-spot table is
not the same thing as interpolation-based convergence analysis.

The domain sensitivity rows show no selected target-region price change for these configurations.
Because \(\Delta S\) and \(N\) are held comparable, this suggests that \(S_{\max}=4K\) is not
distorting the selected target region for these representative cases.

\section{Boundary Stability Results}

\begin{longtable}{P{0.28\textwidth}P{0.15\textwidth}P{0.17\textwidth}P{0.17\textwidth}}
\toprule
Case & Reference & Max selected boundary shift & Interpretation \\
\midrule
\endhead
\path{american_put_grid_80} & grid 180 & 0.0255974 & Coarse boundary visibly differs from reference. \\
\path{american_put_grid_120} & grid 180 & 0.0219485 & Boundary difference remains but is moderate. \\
\path{dividend_call_grid_80} & grid 180 & 0.0220919 & Coarse boundary differs from reference. \\
\path{dividend_call_grid_120} & grid 180 & 0.0217417 & Moderate boundary difference. \\
\path{american_put_domain_smax4} & domain 6 & 0 & Domain extension did not move selected boundaries. \\
\path{dividend_call_domain_smax4} & domain 6 & 0 & Domain extension did not move selected boundaries. \\
\bottomrule
\end{longtable}

Boundary curves remain financially plausible: put boundaries are below the strike, and
dividend-call boundaries are above the strike. Grid-level shifts should not be overread as exact
boundary convergence, because Ticket 09 extraction is threshold-based and nearest-time matching is
used here.

\section{LCP Diagnostic Stability Results}

\begin{longtable}{P{0.29\textwidth}P{0.16\textwidth}P{0.16\textwidth}P{0.17\textwidth}}
\toprule
Case & Obstacle violation & Equation violation & Complementarity max \\
\midrule
\endhead
\path{american_put_grid_80} & 0 & \(1.30\times10^{-9}\) & \(3.48\times10^{-11}\) \\
\path{american_put_grid_180} & 0 & \(1.64\times10^{-9}\) & \(3.25\times10^{-11}\) \\
\path{dividend_call_grid_80} & 0 & \(1.67\times10^{-9}\) & \(5.26\times10^{-11}\) \\
\path{dividend_call_grid_180} & 0 & \(3.58\times10^{-9}\) & \(2.75\times10^{-11}\) \\
\path{american_put_domain_smax6} & 0 & \(1.56\times10^{-9}\) & \(6.65\times10^{-11}\) \\
\path{dividend_call_domain_smax6} & 0 & \(2.14\times10^{-9}\) & \(2.38\times10^{-11}\) \\
\bottomrule
\end{longtable}

All PSOR steps converged across the Ticket 11 runs. LCP diagnostic magnitudes remain near solver
tolerance, and no new obstacle violation appears under refinement or domain extension.

\section{Greek Diagnostic Stability Results}

\begin{longtable}{P{0.29\textwidth}P{0.17\textwidth}P{0.18\textwidth}P{0.15\textwidth}}
\toprule
Case & Full max \(|\Gamma|\) & Strict max \(|\Gamma|\) & Caution \\
\midrule
\endhead
\path{american_put_grid_80} & 20 & 1.25684 & Coarse grid damps Gamma. \\
\path{american_put_grid_180} & 45 & 4.39202 & Finer grid exposes sharper curvature. \\
\path{dividend_call_grid_80} & 20 & 1.64950 & Coarse grid damps Gamma. \\
\path{dividend_call_grid_180} & 45 & 3.83192 & Finer grid exposes sharper curvature. \\
\path{american_put_domain_smax6} & 30 & 3.18002 & Domain-stable with same \(\Delta S\). \\
\path{dividend_call_domain_smax6} & 30 & 2.67989 & Domain-stable with same \(\Delta S\). \\
\bottomrule
\end{longtable}

The Gamma rows are diagnostic, not risk labels. The increasing full-grid maximum is expected
because \(\Gamma\) is a second derivative and sharper grids resolve sharper kink/boundary
features. This confirms Ticket 10's warning: Greek outputs require masking and careful
interpretation.

\section{Figure Interpretation}

\begin{figure}[H]
  \centering
  \includegraphics[width=0.88\textwidth]{results/01_solver_validation/figures/ticket_11_price_vs_grid_size.png}
  \caption{At-the-money selected prices versus grid size for the two representative cases.}
\end{figure}

\begin{figure}[H]
  \centering
  \includegraphics[width=0.88\textwidth]{results/01_solver_validation/figures/ticket_11_price_error_vs_reference.png}
  \caption{Selected price differences relative to the chosen grid/domain reference.}
\end{figure}

\begin{figure}[H]
  \centering
  \includegraphics[width=0.88\textwidth]{results/01_solver_validation/figures/ticket_11_boundary_grid_comparison.png}
  \caption{Approximate boundary curves across grid levels. These are threshold-based diagnostics, not exact analytical boundaries.}
\end{figure}

\begin{figure}[H]
  \centering
  \includegraphics[width=0.88\textwidth]{results/01_solver_validation/figures/ticket_11_domain_cutoff_comparison.png}
  \caption{Domain-cutoff sensitivity is flat in the selected target region under comparable \(\Delta S\).}
\end{figure}

\begin{figure}[H]
  \centering
  \includegraphics[width=0.88\textwidth]{results/01_solver_validation/figures/ticket_11_lcp_diagnostic_stability.png}
  \caption{LCP diagnostic quantities stay near solver tolerance across sensitivity settings.}
\end{figure}

\begin{figure}[H]
  \centering
  \includegraphics[width=0.88\textwidth]{results/01_solver_validation/figures/ticket_11_gamma_stability.png}
  \caption{Gamma diagnostics are grid-sensitive, which is expected for second derivatives near non-smooth features.}
\end{figure}

\section{Testing Strategy and Results}

Ticket 11 tests cover configuration construction, comparable domain spacing, nearest-grid spot
matching, synthetic reference differences, small real solver runs, boundary/LCP/Greek diagnostic
rows, CSV artifacts, figure artifacts, scope exclusions, and Ticket 11 docstrings.

The focused Ticket 11 test command returned:
\begin{lstlisting}
Ran 13 tests in 1.439s
OK
\end{lstlisting}

The full final suite is rerun after report generation and recorded in the chat completion.

\section{Implementation Pitfalls and Debugging Notes}

\begin{longtable}{P{0.25\textwidth}L{0.24\textwidth}L{0.24\textwidth}L{0.16\textwidth}}
\toprule
Event & Symptom & Fix & Lesson \\
\midrule
\endhead
Initial red run & Ticket 11 tests failed because \path{sensitivity.py} and the experiment script did not exist. & Implemented the diagnostics module and experiment. & Red tests should fail for the intended missing capability. \\
Intermediate run & Module tests passed, but experiment setup failed because \path{09_grid_domain_sensitivity.py} was still missing. & Added the experiment script and artifact writers. & Build feature layers in order: API first, then artifact orchestration. \\
\bottomrule
\end{longtable}

Likely traps for future sensitivity work:
\begin{itemize}
  \item confusing grid refinement with domain enlargement;
  \item comparing selected points without reporting actual moneyness;
  \item overclaiming convergence from three grid levels;
  \item treating Gamma outputs as stable labels;
  \item creating stress maps before the validation ladder is reviewed.
\end{itemize}

\section{Limitations}

\begin{longtable}{P{0.32\textwidth}L{0.48\textwidth}}
\toprule
Limitation & Explanation \\
\midrule
\endhead
Only three grid levels & Enough for a preliminary sensitivity gate, not a formal convergence proof. \\
Nearest-node comparison & Selected prices report actual moneyness, but interpolation is not implemented in this ticket. \\
Representative cases only & The study covers one put and one dividend-paying call setup. \\
Gamma remains fragile & Greek sensitivity is diagnostic and should not be exported as labels. \\
No stress maps or datasets & Ticket 11 intentionally stops before risk-surface generation. \\
\bottomrule
\end{longtable}

\section{Lessons Learned}

For greenhand researchers, the main lesson is that numerical validation has layers. A price can
look stable while a boundary or Gamma diagnostic still deserves caution. Ticket 11 shows why the
project records selected spot values, actual moneyness, boundary shifts, LCP diagnostics, and
Greek metrics together.

\section{Link to the Next Ticket}

The next ticket can assemble the solver-validation findings into a broader review gate. Stress
maps, datasets, and neural surrogate work should still wait for human review of the full
validation ladder.

\section{Reproducibility Record}

\begin{longtable}{P{0.35\textwidth}L{0.45\textwidth}}
\toprule
Item & Record \\
\midrule
\endhead
Source & \file{src/american_risk_surfaces/diagnostics/sensitivity.py}. \\
Experiment & \file{experiments/09_grid_domain_sensitivity.py}. \\
Tests & \file{tests/test_grid_domain_sensitivity.py}. \\
Tables & Five \path{ticket_11} CSVs under \file{results/01_solver_validation/tables/}. \\
Figures & Six \path{ticket_11} PNGs under \file{results/01_solver_validation/figures/}. \\
Report & This \LaTeX{} source and compiled PDF. \\
\bottomrule
\end{longtable}

\end{document}
